% !TEX program = xelatex
% MWE for #366: counting-rod numerals in zhnumber
\documentclass[12pt]{article}
\usepackage[margin=1.5cm,paperwidth=13cm,paperheight=15.5cm]{geometry}
\usepackage{fontspec}
\usepackage{zhnumber}

\newfontfamily\RodFont{LXGWWenKaiGBLite-Regular.ttf}
\newfontfamily\JuliaFont{JuliaMono-Regular.ttf}
\setmainfont{LXGWWenKaiGBLite-Regular.ttf}

\zhrodsetup{font=\RodFont}
\pagestyle{empty}

\begin{document}
\obeylines
\textbf{1. 基本用法（可展开）}
\zhrod{12345} \quad \zhrod{-12030.405}

\textbf{2. units：个位取纵式还是横式}
units=vertical\ \ \zhrod{123456789}
{\zhrodsetup{units=horizontal}units=horizontal \zhrod{123456789}}

\textbf{3. zero：零位填〇还是省略}
zero=fill\ \ \zhrod{12030}
{\zhrodsetup{zero=omit}zero=omit\ \zhrod{12030}}

\textbf{4. minus：独立＼还是组合叠加}
minus=slash\ \ \zhrod{-1234}
{\JuliaFont\zhrodsetup{minus=overlay}minus=overlay \zhrod{-1234}}

\textbf{5. \string\zhrodbox：压紧字距（不可展开）}
\zhrod{12345}（未压缩）
\zhrodbox{12345}（已压缩）
\end{document}
